\subsection{A partition}
\label{sec:ha:enabled:partition}

Both instances are running and neither can see the other. Each is healthy, each
believes the other may be dead, and neither can tell which.

This is the condition the arbiter exists for, and it is the only one in this
chapter where doing nothing and doing the obvious thing are both wrong.

\subsubsection{What is at risk}
\label{sec:ha:enabled:partition:risk}

An instance that cannot see its peer has two explanations available and no way to
choose between them. If it promotes itself and the peer is alive, both act and
their states diverge --- and neither can afterwards be shown to have been right.
If it declines and the peer is genuinely dead, nothing acts. That failure is
harder to notice than the first: nothing crashes and nothing logs an error, so it
is found when somebody asks why trading has stopped.

\subsubsection{What the venue must do}
\label{sec:ha:enabled:partition:behaviour}

\begin{requirement}{R-0086}{An instance that cannot see its peer refers the question rather than deciding it}
  Where an instance cannot reach its peer, it shall ask the arbiter which
  instance may act, and shall not decide for itself.
  \rationale{This is the one case an instance cannot settle alone, because it is
  the only one where a second claimant may exist unseen. Where a peer \emph{is}
  visible, the two settle it between themselves --- \reqref{R-0072} --- and the
  two rules are mutually exclusive by construction, so no group has two parties
  entitled to decide.}
  \uncovered{no scenario partitions a pair}
\end{requirement}

\begin{requirement}{R-0087}{A partition does not produce two instances acting}
  A partition of any duration shall not result in both instances of a group
  acting.
  \rationale{\reqref{R-0058} states this generally; it is restated here because
  a partition is the condition that produces it, and because a design can satisfy
  the general statement in every case a test contrives and fail it here.}
  \uncovered{no scenario partitions a pair}
\end{requirement}

\subsubsection{When only one side can see the other}
\label{sec:ha:enabled:partition:asymmetric}

A partition need not be symmetric. One instance may be able to reach its peer
while the peer cannot reach it, and the two rules for settling an entitlement
then both apply at once: the instance that can see its peer settles it locally
under \reqref{R-0072}, while the instance that cannot refers it under
\reqref{R-0086}.

That the two rules are mutually exclusive rests on visibility being the same in
both directions, and here it is not. What closes the case is that an instance
which takes an entitlement reports it immediately rather than at its next
scheduled report, so the party being asked already holds a record of a present
holder and confirms that holder instead of granting the entitlement afresh.

The window is therefore as long as it takes a new holder to report itself, and
not as long as the partition. It is narrowed rather than eliminated, which is a
weaker guarantee than one party alone issuing entitlements would give --- and it
is the case most worth producing deliberately, because reasoning is a poor
substitute for having seen it.

\subsubsection{What this costs, stated plainly}
\label{sec:ha:enabled:partition:cost}

An instance on the losing side of a partition is not stopped. It keeps running,
keeps believing what it believed, and is refused at every interaction it attempts
--- \reqref{R-0065}. It cannot corrupt anything, because nothing it sends is
accepted, but it also does not stop, and it may hold resources until somebody
intervenes.

That is the accepted trade. Removing the instance instead would require
authority over the machine it runs on: a separate management network, hardware
power control, and credentials to use them --- properties of a deployment rather
than of this software, absent from a developer's machine, and therefore a
mechanism that could only ever be exercised in one environment. Refusing
everything a stale instance sends needs none of that and is tested wherever the
venue runs.
